\documentclass[10pt,a4paper]{amsart}
\usepackage[margin=19mm]{geometry}
\usepackage{amsmath,amssymb,mathtools}
\usepackage[scr=rsfs,cal=euler]{mathalfa}
\usepackage{xfrac}
\usepackage[unicode,hidelinks,bookmarks=false]{hyperref}
\newcommand{\ie}{i.e.\ }
\newcommand{\cf}{cf.\ }
\newcommand{\resp}{respectively\ }
\newcommand{\Subsection}{\subsection}
\providecommand{\nobreakdash}{\nobreak\hbox{-}}
\setlength{\emergencystretch}{2em}
\multlinegap0pt
\usepackage{bm}
\usepackage{bbm}
\usepackage{dsfont}
\usepackage{xspace}
\usepackage{cancel}
\usepackage{mathtools}
\usepackage{amsthm}
\makeatletter
\@ifundefined{thm}{\newtheorem{thm}{Theorem}}{}
\@ifundefined{defi}{\newtheorem{defi}[thm]{Definition}}{}
\@ifundefined{lem}{\newtheorem{lem}[thm]{Lemma}}{}
\@ifundefined{prop}{\newtheorem{prop}[thm]{Proposition}}{}
\makeatother
\def\bar{\overline}
\newcommand{\HH}{\mathcal{H}}     % for the set of Halfspaces or the functor
\newcommand{\NN}{\mathbb{N}}      % for natural numbers
\newcommand{\YY}{\mathcal{Y}}     % for the set of measurable section to $Y'$
\newcommand{\hh}{\mathfrak{h}}    % for denoting halfspaces
\title[Cocycle superrigidity for median spaces of finite rank]{Cocycle superrigidity for median spaces of finite rank}
\author{Biao Ma, Lamine Messaci}
\date{Source: arXiv:2506.19043v2 (CC BY 4.0)}
\begin{document}
\maketitle

\noindent\emph{Source and licence.} Cocycle superrigidity for median spaces of finite rank. Version arXiv:2506.19043v2 (CC BY 4.0), distributed under the Creative Commons Attribution 4.0 International licence, as recorded in the arXiv licence metadata for this exact version and shown on the abstract page https://arxiv.org/abs/2506.19043v2. Full rights notice: licenses/arxiv-2506.19043-CC-BY-4.0.txt.\par\smallskip

\noindent\textit{Editorial scope.} Counted unit: the Thm whose source label is \texttt{theoerem\_superrigidity\_cocycle} in the article (source0.tex lines 1464--1466, source label \texttt{theoerem\_superrigidity\_cocycle}), delivered together with the article's own complete original proof of it (source0.tex lines 1467--1497, one \texttt{proof} environment of 31 lines). No mathematics is written, repaired or re-derived: statement and proof are the article's own text line for line. Required context retained uncounted: def\_median\_core (lines 772--780); thm\_countable\_medians\_subspace (lines 844--846); lem\_Fubini2 (lines 1375--1381); prop\_boundary\_map\_factors (lines 1447--1449). Every definition, display and label of those lines is retained unchanged. Omitted from this package and not redistributed: 1--771, 781--843, 847--1374, 1382--1446, 1450--1463, 1498--1622. Nothing inside a retained excerpt is omitted or altered: every retained body is delivered exactly as the article's lines are written, so no omission receipt of any kind accompanies this package. Citations: the retained excerpts carry no citation command at all, so no bibliography entry of the article is transcribed and no reference is redistributed. Reference adaptation: none was needed and none was made; every reference inside the delivered unit resolves inside it, so no unresolved reference or citation is produced. Printed slips kept literal: no printed word, symbol, hypothesis or number of the counted statement or of its proof was altered. Layout and class: the article's own class is \texttt{amsart} with packages including amsmath, xcolor, graphics, amsfonts, epsfig, float, graphicx, epstopdf, amssymb, mathrsfs, amsthm, hyperref, cleveref, natbib; the wrapper is the lane's pinned \texttt{amsart} editorial wrapper at 10pt on A4, which restates the theorem-like environments the retained text needs and adds the article's own macro definitions that the retained text needs (\texttt{\textbackslash HH}, \texttt{\textbackslash NN}, \texttt{\textbackslash YY}, \texttt{\textbackslash bar}, \texttt{\textbackslash hh}), with extra wrapper packages (bm, bbm, dsfont, xspace, cancel, mathtools). The delivered numbering is the unit's own. Assets: no retained span contains a figure environment, a graphics-inclusion command, a PDF-page inclusion, a table or a TikZ picture; no image file is copied into this package, the package declares no asset dependency and no image file is redistributed with it. Licence: version arXiv:2506.19043v2 of this article is distributed under the Creative Commons Attribution 4.0 International licence, as recorded in the arXiv licence metadata for this exact version and shown on the abstract page https://arxiv.org/abs/2506.19043v2. A full rights notice is delivered at licenses/arxiv-2506.19043-CC-BY-4.0.txt. Characters: every retained excerpt is pure ASCII, so the delivered engine, XeLaTeX with its default Latin Modern text font, reports no missing character.
\begin{defi}[\textbf{\textit{Median core}}]\label{def_median_core}
    Let $X$ be a complete irreducible separable median space of finite rank. 
    Consider the following $Isom(X)$-invariant subset
  \[
    Y_0:=\{m_{\bar X}(\tilde{x}_1,\tilde{x}_2,\tilde{x}_3) \ |\ \tilde{x}_1,\tilde{x}_2,\tilde{x}_3\in  \partial_r X, \ x_i\neq x_j\}\subseteq X 
  \]
    and we define recursively $Y_{i+1}:= m(Y_i,Y_i,Y_i)$. We set $Y':=\displaystyle{\bigcup_{i\in\NN}Y_i}$ and remark that it is a $Isom(X)$-invariant median subspace.\par
   We call the closure of $Y'$ in $X$ the $\textbf{\textit{median core}}$ and denote the closure by $Y$.% be $Y:=\bar{Y'}\cap X$, where $Y'$ is the median space arising in the proof of Theorem \ref{thm_countable_medians_subspace}.
\end{defi}

\begin{thm}\label{thm_countable_medians_subspace}
    The median subspace $Y'$ arising in Definition \ref{def_median_core} is a countable $Isom(X)$-invariant median subspace.
\end{thm}

\begin{lem}\label{lem_Fubini2}
    For any fixed thick halfspace $\hh\in \HH(Y)$ and a fixed $\epsilon>0$. Then for almost all $\omega \in\Omega$ the following 
    \[
    B_{\hh,\omega}^\varepsilon:=\{b\in B\ | \ ev_\hh(\Phi(\omega,b))>1-\varepsilon\}
    \]
    is of positive measure.
\end{lem}

\begin{prop}\label{prop_boundary_map_factors}
    Let $G:=G_1\times G_2$ where $G_1$ and $G_2$ are locally compact, second countable, non compact groups. Then the boundary maps $\phi$ factors through the projection onto $\Omega\times B_i$ for one of the $B_i$'s.
\end{prop}

\begin{thm}\label{theoerem_superrigidity_cocycle}
    There exists a group homomorphism $h:G\rightarrow Isom(Y)$ that factors through one of the $G_i$'s, and $h$ is cohomologous to the cocycle $\alpha$
\end{thm}

\begin{proof}
Let us denote by $Y_0\subseteq Y$ the $\alpha$-invariant subset obtained by taking the median points of distinct triple in $\partial_r X\subseteq \bar Y$ and let us denote by $Y'$ its median hull inside $X$ as in Definition \ref{def_median_core}. The median subspace $Y'$ is countable and $\alpha$-invariant (see Theorem \ref{thm_countable_medians_subspace}). \par 
 By Proposition \ref{prop_boundary_map_factors}, the boundary map $\phi$ factors through $\Omega\times B_i$ for one of the $B_i$'s. Let us assume that it is $\Omega\times B_1$. Then almost all fixed $b_1$ gives rise to an $\alpha_{G_2}$ equivariant section, where $\alpha_{G_2} = \alpha|_{G_2}$ is the cocycle on $G_2$ obtained as the restriction of $\alpha$ to $\{e_{G_1}\}\times G_2$. We use the latter familly to construct infinitely many $\alpha_{G_2}$-equivariant section from. Let $\YY$ be the set of equivalence classes of measurable $\alpha_{G_2}$-invariant sections from $\Omega$ to $Y'$, where two measurable section are equivalent if they coincide in a subset of total measure.\par 
 \textit{\textbf{Claim 1:} The set $\YY$ is not empty. }\par 
 Let us fix $y\in Y_0$ and consider $\tilde{x}_1,\tilde{x}_2,\tilde{x}_3\in\partial_r X$ such that $y=m(\tilde{x}_1,\tilde{x}_2,\tilde{x}_3)$. There exists then three strongly separated thick halfspaces $\hh_1,\hh_2,\hh_3\in\HH(X)$ such that $\tilde{x}_i\in\tilde{\hh}_i$ and for any $\tilde{x}_i'\in\tilde{\hh}_i$ we have $m(\tilde{x}_1,\tilde{x}_2,\tilde{x}_3)=y$. Then for any measurable subset $A\subseteq \Omega$ of positive measure, Lemma \ref{lem_Fubini2} and Fubini's Theorem ensures the existence of a measurable subset $A'\subseteq A$ of positive measure  and $b_1,b_2,b_3\in B_1$ such that for any $\omega\in A'$, we have $\phi(\omega,b_i)\in\tilde{\hh}_i$. The measurable decomposition $m(\partial_r X):=\partial_r X\sqcup Y_0$ is $Isom(X)$-invariant. Hence, by the $\alpha_{G_2}$-equivariance of the sections and the ergodicity of the $G_2$ action on $\Omega$, we deduce that for almost all $b_1,b_2,b_3\in B$ we have $m(\phi(\omega,b_1),\phi(\omega,b_2),\phi(\omega,b_3))\in Y_0$, which gives a well defined class of $\alpha_{G_2}$-invariant measurable section from $\Omega$ to $Y_0$ which completes the first claim.  
% \begin{eqnarray*}
%    \phi_{b_1,b_2,b_3}:\Omega &\rightarrow & Y_0\\
%\ \ \ \ \omega &\mapsto & m(\phi(\omega,b_1),\phi(\omega,b_2),\phi(\omega,b_3))
 %   \end{eqnarray*}
 \par
 The action of $G_2$ on $\Omega$ is ergodic and the target space $Y'$ is countable, the set $\YY$ is also countable. Hence let us fix a choice of representative of each class of $\YY$ to obtains a sequence $(f_i)_{i\in\NN}$, that is $\YY:=\displaystyle{\bigcup_{i\in\NN}[f_i]}$, where $[f_i]$ is the equivalent class of $f_i$.\par 
  \textit{\textbf{Claim 2:} Almost everywhere, the evaluation map give rise to a surjection between $\YY$ and $Y'$.}\par
  In the proof of the first claim, we have shown that for any measurable subset $A\subseteq \Omega$ of positive measure and any $y\in Y'$, there exists a measurable subset of positive measure $A'\subseteq A$ and $f\in \YY$ such that $f(\omega)=y$ for any $\omega\in A'$ (note that we did it when $y$ lies in $Y_0$ but the same argument applies for any point in $Y'$ by iterating the median operator on the $\alpha_{G_2}$ invariant sections obtained in the previous step and taking value in $Y_0$). Hence, the measurable subset 
  \[
  A_y:=\bigcup_{i\in\NN}f_i^{-1}(y)
  \]
 is of total measure, as it is measurable and it intersects any measurable subset of positive measure in $\Omega$.\par 
 After setting $\Omega':=\displaystyle{\bigcap_{y\in Y'}A_y}$, we deduce that for each $\omega\in \Omega'$ the evaluation map $ev_\omega:\YY\rightarrow Y'$ is surjective.\par 
 By the ergodicity of the $G_2$-action on $\Omega$ and the equivariant of the sections $f_i's$, $\YY$ is naturally endowed with a median metric $d_\YY(f_i,f_j):=\int_{\Omega}d_Y(f_i(\omega),f_j(\omega))d\mu(\omega)$. \par 
 
\textit{\textbf{Claim 3:} Almost everywhere, the evaluation map give rise to an isometric identification between $\YY$ and $Y'$.}\par 
For any $i,j\in\NN$, the function $d(f_i(*),f_j(*))$ is $G_2$-invariant, hence it is constant by the metric ergodicity of the $G_2$-action on $\Omega$. Thus, there exists a subset $\Omega''\subseteq \Omega$ of total measure, such that $f_i(\omega)\neq f_j(\omega)$ for any $\omega\in\Omega''$ and any distinct $i,j\in\NN$ (the family $(f_i)_{i\in\NN}$ is assumed to be pairwise essentially distinct). Therefore, the evaluation map $ev_\omega:(\YY,d_\YY)\rightarrow (Y',d_Y)$ is an isometric bijection for every $\omega\in \Omega''$.\par 

After fixing a base point $\omega_0$ from which we identify $\YY$ with $Y'$, we have a measurable map 
\begin{eqnarray*}
    F:\Omega'' \times Y' &\rightarrow & Y'\\
\ \ \ \ (\omega,f_i(\omega_0)) &\mapsto & f_i(\omega)
    \end{eqnarray*}
Hence, each fixed $\omega\in\Omega''$ gives rise to an element in $Isom(Y')$. Thus, we have a measurable map $h:\Omega\rightarrow Isom(Y')\cong Isom(Y)$.\par 
Therefore, the $G_1$-isometric action on $\YY$ is cohomologous to the cocycle $\alpha$ under conjugation by $h$.
\end{proof}

\end{document}
